% translate_docs: manuals/preconditioning_manual.tex @ 2a57ad48e16e3e1949a46e871e394e3d2231476f
\documentclass[superscriptaddress,onecolumn,floatfix]{revtex4-2}

\usepackage{graphicx}
\usepackage{epsfig}
\usepackage{color}
\usepackage{amssymb}
\usepackage{amsmath}
\usepackage{array}
\usepackage[loose]{units}
\usepackage[russian]{babel}
\usepackage[unicode]{hyperref}
\usepackage{cleveref}
\usepackage{tasks}
\usepackage[caption=false]{subfig}
\usepackage{letltxmacro}

\LetLtxMacro{\ORIGselectlanguage}{\selectlanguage}
\makeatletter
\DeclareRobustCommand{\selectlanguage}[1]{%
    \@ifundefined{alias@\string#1}
      {\ORIGselectlanguage{#1}}
      {\begingroup\edef\x{\endgroup
         \noexpand\ORIGselectlanguage{\@nameuse{alias@#1}}}\x}%
}
\newcommand{\definelanguagealias}[2]{%
  \@namedef{alias@#1}{#2}%
}
\makeatother

\definelanguagealias{ru}{russian}


\begin{document}
\selectlanguage{ru}

\begin{abstract}
  В данном руководстве описаны основные принципы использования и ответы на часто задаваемые вопросы, касающиеся комплекта для ручной предварительной подготовки оборудования.
\end{abstract}
\title{Инструкция по комплекту для предварительной подготовки}
\author{Бутик электронных кнопок}
\email{support@electroniqbuttons.com}
\date{\today}

\maketitle


\section{Какие автомобили могут это использовать?} \label{sect:eligibility}
Данный комплект предназначен для работы с большинством моделей Ioniq 5s 2021-2024 годов выпуска, Ioniq 6s 2022-2025 годов выпуска и EV6s 2022-2024 годов выпуска. Мы протестировали его на Ioniq 5s из трех разных стран, а также на одном Ioniq 6 и одном EV6. \href{mailto:requests@electroniqbuttons.com}{Связаться с нами} Если вас интересует использование комплекта на другом автомобиле, имейте в виду, что он будет работать только при наличии подогревателя батареи. Для этого проверьте наличие меню «Режим кондиционирования батареи» в разделе «Электромобиль». $\to$ Настройки электромобиля (см. рис.) \ref{subfig:menus:EVsettings}) на головном устройстве вашего автомобиля (центральный экран) (см. рис. \ref{subfig:menus:batt_cond}).
\begin{figure}[h]
  \subfloat[]{\label{subfig:menus:EVsettings}
  \includegraphics[height=1.95in]{photos/EV_settings.jpg} }
  \subfloat[]{\label{subfig:menus:batt_cond}
  \includegraphics[height=1.95in]{photos/battery_conditioning_mode.jpg} }
   \caption{(a) Как найти меню настроек EV: на главном экране нажмите значок «EV»; затем в появившемся меню слева нажмите «Настройки EV». (b) Экран «Режим восстановления батареи». \label{fig:menus}}
\end{figure}

Если у вас нет этого экрана, вам не следует покупать этот комплект. Пожалуйста, свяжитесь с дилером или знающими автовладельцами в вашей стране, чтобы оценить возможность установки этой функции на ваш автомобиль с дополнительными действиями. Для США, Канады или Норвегии, пожалуйста, свяжитесь с нами по адресу [адрес электронной почты]. \href{mailto:support@electroniqbuttons.com}{support@electroniqbuttons.com} И мы с удовольствием поможем вам оценить это, предоставив немного информации о вашем автомобиле.

Если вместо «Режима восстановления батареи», как показано на рисунке, используется другой режим. \ref{subfig:menus:batt_cond}Если у вас включен «Зимний режим», вашему автомобилю просто необходимо обновить прошивку в дилерском центре. Найдите соответствующий технический бюллетень для вашей страны, касающийся этого обновления, и обратитесь в ближайший компетентный дилерский центр. После выполнения этого шага и подтверждения того, что меню «Режим восстановления батареи» заменяет «Зимний режим», вы можете приобрести этот комплект.

\section{Базовая комплектация}
Начните с \href{https://github.com/tylerharvey/Ioniq5_CAN/blob/main/guides/}{инструкции по установке}. Видеть \href{https://meatpihq.github.io/wican-fw/config/wifi}{Документация WiCAN} Для базовой настройки WiCAN. Активация предварительной подготовки не обязательна, но настоятельно рекомендуется хотя бы изменить пароль Wi-Fi по умолчанию в целях безопасности.  

\section{Пользовательский интерфейс}
Мы активно разрабатываем пользовательский интерфейс, поэтому заглядывайте сюда почаще!

Для настройки режима ручной предварительной подготовки подключитесь к сети WiCAN WiFi и перейдите на страницу конфигурации WiCAN по адресу: \href{http://192.168.80.1}{http://192.168.80.1}Третья вкладка, называемая «Кнопки», показана на рисунке. \ref{subfig:ui:buttons} Позволяет настроить кнопку ручной предварительной подготовки, а также любые другие пользовательские кнопки, которые будут добавлены позже.

\begin{figure}[h]
  \subfloat[]{\label{subfig:ui:buttons}
  \includegraphics[height=1.95in]{photos/ui_2026-07-04.png} }  \\
   \caption{Скриншот страницы «Кнопки» в конфигурации WiCAN HTTP.}
\end{figure}

Первый вариант, «Режим предварительной подготовки», в настоящее время находится в разработке и имеет только один допустимый режим: один раз. После запуска одного цикла предварительной подготовки она прекратится. Мы работаем над режимами «Непрерывный» и «Постоянный», которые позволят системе управления батареей (BMS) автоматически повторно нагревать батарею, если температура падает, либо до выключения автомобиля («Непрерывный»), либо даже во время циклов включения/выключения автомобиля («Постоянный»).

Второй вариант позволяет выбрать кнопку для назначения предварительной подготовки. Дополнительные параметры будут доступны в режиме "человек посередине" с помощью настраиваемой системы WiCAN. По умолчанию отображаются 5 вариантов:
\begin{enumerate}
  \item Неполноценный
  \item Кнопка «Звезда» (на рулевом колесе) — по умолчанию
  \item Кнопка «Звезда» (Климатическая шкала)
  \item Кнопка тюнера (Климат-бар)
  \item Регулятор громкости (климатическая шкала)
\end{enumerate}
При выборе значения «Отключено» ни одна кнопка не будет запускать предварительную подготовку. Это может быть желательно при обращении в автосервис, чтобы избежать путаницы у специалистов. Варианты «нажать» означают нажатие этих кнопок, а не поворот или нажатие вверх или вниз. Если ни один из этих вариантов не работает, вы можете установить флажок «Показать все параметры кнопок», чтобы включить дополнительные параметры:
\begin{tasks}(2)
  \task Режим (рулевое колесо)
  \task Говори (руль)
  \task Вызов (рулевое колесо)
  \task Регулировка громкости (на рулевом колесе)
  \task Увеличение громкости (на рулевом колесе)
  \task Уменьшение громкости (на рулевом колесе)
  \task Перемотать вверх (руль)
  \task Перемотать вниз (руль)
  \task ОК (руль)
  \task Карта (Климатическая шкала)
  \task Навигация (Климат-панель)
  \task Медиа (Климат-бар)
  \task Настройка (климат-бар)
  \task Снижение тональности (Климатическая полоса)
\end{tasks}
Выберите кнопку, которая вам подходит! Мы надеемся вскоре предложить и физические значки.

Для ручной активации предварительной подготовки просто нажмите выбранную кнопку активации. Сразу же на приборной панели появится обратный отсчет в метрах/ярдах, который перехватывает управление индикатором расстояния до пункта назначения. Если у вас включен круиз-контроль, он блокирует наиболее заметный индикатор расстояния до пункта назначения. Используйте кнопку «Страницы» на рулевом колесе, чтобы перейти на страницу приборной панели с навигационными инструкциями. Предварительная подготовка начинается, когда этот обратный отсчет достигнет нуля. Если обратный отсчет обнуляется, это означает, что произошло что-то непредвиденное, и попытка повторяется. 

Когда активируется предварительная подготовка, появится стандартное всплывающее окно с сообщением «Предварительная подготовка батареи активирована для оптимальной зарядки постоянным током», как показано на рисунке. \ref{subfig:precon:popup}Кроме того, над символом батареи в левом нижнем углу приборной панели появится значок снежинки или красной катушки, как показано на рисунке. \ref{subfig:precon:coil}.

\begin{figure}[h]
  \subfloat[]{\label{subfig:precon:popup}
  \includegraphics[height=1.95in]{photos/precon_active.jpg} }  \\
  \subfloat[]{\label{subfig:precon:coil}
  \includegraphics[height=1in]{photos/coil_on_battery.jpg} }
   \caption{(a) Скриншот всплывающего окна, указывающего на успешную активацию предварительной подготовки батареи. (b) Пример скриншота значка катушки, наложенного на символ батареи в левом нижнем углу приборной панели. \label{fig:precon}}
\end{figure}

Если предварительная подготовка так и не активируется, убедитесь, что выполнены встроенные условия для ее активации:
\begin{enumerate}
  \item Температура батареи должна быть ниже $\unit[21]{C}$/$\unit[70]{F}$
  \item Уровень заряда должен быть выше $\unit[20]{\%}$
  \item Экран «Режим восстановления батареи» присутствует в меню «Настройки электромобиля» (см. раздел). \ref{sect:eligibility})
\end{enumerate}
К сожалению, это ограничения, заложенные в прошивке блока управления батареей или в аппаратной части автомобиля, и мы не можем изменить их с помощью этого комплекта.

\section{Потребление электроэнергии}
Этот комплект работает от 12-вольтовой шины вашего автомобиля, а это значит, что он может потреблять энергию от 12-вольтового аккумулятора вашего автомобиля. Стоит обсудить несколько моментов, связанных с этим:
\begin{enumerate}
  \item Сколько энергии он потребляет?
  \item Не повредит ли это 12-вольтовому источнику питания?
  \item MITM против параллельного соединения
\end{enumerate}

\subsection{Сколько энергии он потребляет?}

Мы измерили энергопотребление при выключенном двигателе (12,7 В) для нескольких микроконтроллеров в различных состояниях. Тесты WiCAN в режиме «обычного использования» проводились без каких-либо HTTP-запросов и без подключения BLE через CarScanner или аналогичные программы. Эти тесты проводились на выбираемом CAN-адаптере, где, насколько нам известно, в настоящее время невозможны запросы OBD, поэтому микроконтроллеры, поддерживающие стандартный OBD (WiCAN, vLinker MC+), вероятно, потребляли меньше энергии, чем при активном опросе параметров OBD через обычный порт OBD. 

\subsubsection{Оригинальная версия WiCAN 4.20.\_eb0.1.1-beta19}
\begin{enumerate}
  \item{Обычное использование:} 30 мА
  \item{Использование HTTP-страницы:} до 34 мА
  \item{Загрузка:} до 36 мА
  \item{CAN-монитор:} до 38 мА
  \item{Спать:} $<$1 мА
\end{enumerate}

\subsubsection{WiCAN, первая версия, работает под управлением v4.20\_eb0.1.1-beta19}
\begin{enumerate}
  \item{Обычное использование MITM:} 40 мА
  \item{Обычное параллельное использование:} 34 мА
  \item{Загрузка:} до 48 мА
\end{enumerate}

\subsubsection{WiCAN Pro работает под управлением версии 4.50.}
\begin{enumerate}
  \item{Обычное использование:} 50 мА
\end{enumerate}

\subsubsection{Белая панда с запятой, бегущий последний аниматроник\_панда сборка}
\begin{enumerate}
  \item{Обычное использование:} 32 мА
\end{enumerate}

\subsubsection{vLinker MC+}
\begin{enumerate}
  \item{Обычное использование (без подключения):} 15-18 мА
  \item{CarScanner подключен через BLE:} 20-25 мА, пиковые значения до 60 мА
\end{enumerate}

\subsection{Не повредит ли это 12-вольтовому источнику питания?}

В WiCAN доступен режим энергосбережения, называемый спящим режимом, в котором практически все отключается, и устройство потребляет очень мало энергии. Обычно он настраивается на запуск в зависимости от напряжения 12 В в вашем автомобиле. Если вас беспокоит напряжение 12 В, мы рекомендуем включить эту функцию, которую можно найти в интерфейсе WiCAN по адресу: \href{http://192.168.80.1}{http://192.168.80.1} на вкладке «Энергосбережение». Дополнительную информацию см. здесь: \href{https://meatpihq.github.io/wican-fw/config/sleep-mode}{https://meatpihq.github.io/wican-fw/config/sleep-mode}Мы активно работаем над улучшением этой функции, чтобы она работала на основе CAN-шины (то есть, считывала сигнал выключения автомобиля), а не на основе напряжения. 

За 7 дней работы со свежим 12-вольтовым аккумулятором емкостью 60 Ач типичное энергопотребление модифицированной системы WiCAN, составляющее 40 мА, составило бы 40 мА. $40 \unit[\mathrm{mA}] * 24 \unit[\mathrm{hours/day}] * 7 \unit[\mathrm{days}] = 6.7 Ah$, что составляет примерно 10\% of the 12V state of charge. Unlocked and off, a 2022 Canadian Preferred AWD Ioniq 5 with no aftermarket loads draws about 1.75 A (1750 mA). After ECUs quiet down 5 minutes after locking the car, the draw is still 0.75 A (750 mA). (The car has a further "sleep" mode after 7 days of inactivity; this draw hasn't been measured here.) So, a customized WiCAN in MITM mode and not in sleep mode increases draw over the first 7 days of inactivity by 5\%.

\subsection{MITM против параллельного выполнения}
Если жгут проводов комплекта предварительной подготовки находится в режиме MITM (Man-in-the-Middle) и ваш WiCAN переходит в спящий режим, запросы к приложениям Hyundai/Kia не будут выполняться. В наших планах на будущее — добавить частичный спящий режим, который позволит пересылать данные по CAN, но отключит Wi-Fi. Мы рекомендуем оставлять ваш WiCAN в спящем режиме с очень низким напряжением, например, 12,3 В, если жгут проводов находится в режиме MITM и вы хотите использовать запросы к приложениям Hyundai/Kia. 

\section{Часто задаваемые вопросы}
\begin{enumerate}
  \item[В:] Как мне это установить?
  \item[А:] См. \href{https://github.com/tylerharvey/Ioniq5_CAN/blob/main/guides/}{инструкции по установке}Если это кажется слишком сложным, обратитесь за помощью в ближайший магазин товаров для подростков или магазин автозвука!
  \item[В:] Почему после установки перестали работать камера заднего вида, автомобильное приложение или другие функции центрального экрана?
  \item[А:] Проверьте \href{https://github.com/tylerharvey/Ioniq5_CAN/blob/main/guides/harnesses/head_unit_MITM/MITM_harness_modes.pdf}{руководство по режимам крепления ремня безопасности}Если к порту OBD подключено оригинальное устройство WiCAN или нет, жгут проводов должен работать в параллельном режиме. Если у вас модифицированное устройство WiCAN, жгут проводов должен работать в режиме MITM (атака посредника).
  \item[В:] Это облегчит угон моей машины?
  \item[А:] Маловероятно, но конструкция устройства спроектирована таким образом, чтобы его было трудно заметить по нескольким причинам, включая безопасность. Мы настоятельно рекомендуем сменить пароль WiCAN. 
  \item[В:] Это аннулирует мою гарантию?
  \item[А:] Проверьте законодательство вашей страны о гарантиях и защите прав потребителей. В США Закон Магнуссона-Мосса о гарантиях обеспечивает надежную защиту потребителей. Угрозы автопроизводителей о том, что неоригинальные запчасти «аннулируют гарантии», во многих случаях не выдерживают критики в суде. (Это не юридическая консультация.) В то же время диагностика проблем автомехаником представляет собой сложную задачу, и мы настоятельно рекомендуем демонтировать неоригинальные запчасти и аксессуары, которые могут затруднить диагностику. Например, если у вас возникли проблемы с головным устройством или USB-портом, целесообразно демонтировать этот комплект перед обращением в дилерский центр, чтобы избежать осложнений при диагностике. 
  \item[В:] Какова рекомендуемая конфигурация комплекта для проведения базового технического обслуживания моего автомобиля в дилерском центре?
  \item[А:] Мы рекомендуем отключить WiCAN и убедиться, что кабель питания подключен. \href{https://github.com/tylerharvey/Ioniq5_CAN/blob/main/guides/harnesses/head_unit_MITM/MITM_harness_modes.pdf}{параллельный режим}.
  \item[В:] А что, если я захочу построить это сам?
  \item[А:] Сделайте это! Всё, что мы сделали, — \href{https://github.com/tylerharvey/Ioniq5_CAN}{открытый исходный код}Вы можете использовать схему проводки для самостоятельной сборки. Разъемы доступны на AliExpress. Однако, если вы внесете свой вклад в общедоступные знания в виде запроса на изменение (pull request), \href{https://github.com/dragz/egmpdbc}{файл DBC} Или же доступны скидки, которые позволят снизить чистую стоимость до уровня, конкурентоспособного с ценами для самостоятельного ремонта, при значительно меньших затратах на рабочую силу.
  \item[В:] Почему это нельзя сделать с помощью обновления программного обеспечения?
  \item[А:] Вполне возможно. У нас есть свои предположения, почему Hyundai/Kia не включили это в обновление навигационной системы. Если вы знаете, как это сделать. \href{https://discord.gg/zyrngQ5qQ}{взломать вашу головную магнитолу} Если вы нашли функции интерфейса CAN, можете добавить их. \href{https://www.reddit.com/r/KiaEV6/comments/1jdd56u/comment/mid9xie/}{кнопка ручного предварительного прогрева головного устройства}Однако распространение взломанной прошивки в настоящее время невозможно, поэтому сделать широко доступным «обновление» от стороннего производителя будет сложно.
  \item[В:] Могу ли я реализовать это на своем устройстве Comma?
  \item[А:] Мы это протестировали и считаем, что это не работает. По всей видимости, не происходит пересылки сообщений предварительной подготовки с шин CAN A или E, на которых подключен Comma, на шины M или G, на которых подключен BMU.
  \item[В:] Это убьет мое отделение интенсивной терапии?
  \item[А:] Нет.
  \item[В:] Могу ли я использовать встроенную навигацию при работе с этим комплектом?
  \item[А:] В бета-версии это невозможно. Эта проблема будет исправлена в модифицированной версии WiCAN.
  \item[В:] Повлияет ли обновление навигационной системы на возможность использования комплекта?
  \item[А:] Это крайне маловероятно, поскольку комплект не зависит от связи с головным устройством. Обновление блока управления батареей может помочь (например, обновление, которое позволило ретроактивно включить предварительный прогрев на ранних моделях Ioniq 5), но мы считаем, что это должны делать дилерские центры. Однако мы не можем гарантировать это на 100%.\% certainty that a navigation update will not affect the functionality of the kit. In the off chance that this does happen, let us know and we will do our best to update logic to restore manual preconditioning functionality.
  \item[В:] Зачем этот обратный отсчет? Почему нельзя начать предварительную подготовку сразу после нажатия кнопки?
  \item[А:] Блок управления батареей (ЭБУ внутри батареи, содержащий BMS) имеет встроенный 60-секундный тайм-аут после первоначального запроса на предварительную подготовку. Если запрос прекратится до истечения 60-секундного тайм-аута, предварительная подготовка не начнётся. Поэтому мы постарались сделать этот 60-секундный период ожидания очевидным и понятным для пользователя.
  \item[В:] Вы сделаете мне новую кнопку?
  \item[А:] Возможно! Напишите нам по электронной почте: \href{mailto:requests@electroniqbuttons.com}{requests@electronicbuttons.com} с учетом ваших запросов.
  \item[В:] Ваш магазин называется «Electroniq Buttons Boutique». Где находятся сами кнопки?
  \item[А:] Мы работаем над тем, чтобы сделать физические кнопки доступными. У Роя есть протестированный прототип, показанный в его видео. \href{https://github.com/dragz/ironiq}{Ironiq} репозиторий.
\end{enumerate}
% \begin{figure}[h]
%   \subfloat[]{\label{subfig:overview:harness}
%   \includegraphics[height=1.95in]{photos/whole_harness.jpg} }
%   \subfloat[]{\label{subfig:overview:shunt_cap_labeled}
%   \includegraphics[height=1.95in]{photos/labeled_connectors_MITM_harness.jpg} }
%    \caption{(a) Man-in-the-middle harness. (b) Overview of connectors showing shunt cap. \label{fig:overview}}
% \end{figure}
\appendix 
\bibliographystyle{apsrev4-1}
\bibliography{stemh_model}{}
\end{document}